\documentclass{article}
\usepackage[T1]{fontenc}
\usepackage[utf8]{inputenc}
\usepackage{longtable}
\usepackage{geometry}
\usepackage{tikz}
\usepackage{longtable}
\usepackage{polski}
\usetikzlibrary{shapes.geometric, arrows.meta, positioning, calc}

\geometry{margin=2.5cm}

\title{Model skarbca danych - Uczelnia (Data Vault)}
\author{Mikołaj Kubś}
\date{}

\begin{document}

\maketitle

\section{Model wygenerowany przez sztuczną inteligencję - ChatGPT}

\subsection{Dziedzina}

Uczelnia wyższa - model dotyczy danych operacyjnych związanych ze studentami, prowadzącymi oraz kursami.

\subsection{Huby (klucze biznesowe)}

\subsubsection*{HUB\_Student}
\begin{itemize}
  \item \textbf{Student\_ID} - np. numer indeksu
  \item Atrybuty biznesowe: brak (wszystkie zmienne w satelitach)
\end{itemize}

\subsubsection*{HUB\_Prowadzacy}
\begin{itemize}
  \item \textbf{Prowadzacy\_ID} - np. PESEL lub ID uczelniane
  \item Atrybuty biznesowe: brak
\end{itemize}

\subsubsection*{HUB\_Kurs}
\begin{itemize}
  \item \textbf{Kurs\_ID} - np. kod kursu
  \item Atrybuty biznesowe: brak
\end{itemize}

\subsection{Linki (relacje biznesowe)}

\subsubsection*{LINK\_Student\_Kurs}
\begin{itemize}
  \item Student\_ID
  \item Kurs\_ID
  \item Relacja: student zapisany na kurs
\end{itemize}

\subsubsection*{LINK\_Kurs\_Prowadzacy}
\begin{itemize}
  \item Kurs\_ID
  \item Prowadzacy\_ID
  \item Relacja: prowadzący przypisany do kursu
\end{itemize}

\subsubsection*{LINK\_Egzamin}
\begin{itemize}
  \item Student\_ID
  \item Kurs\_ID
  \item Data\_Egzaminu
  \item Relacja: podejście studenta do egzaminu
\end{itemize}

\subsection{Satelity (miary i wymiary)}

\subsubsection*{SAT\_Student\_Dane}
\begin{itemize}
  \item Student\_ID, Imię, Nazwisko, Email
  \item Data\_Wejscia, Data\_Wyjscia
\end{itemize}

\subsubsection*{SAT\_Prowadzacy\_Dane}
\begin{itemize}
  \item Prowadzacy\_ID, Imię, Nazwisko, Tytuł, Email
  \item Data\_Wejscia, Data\_Wyjscia
\end{itemize}

\subsubsection*{SAT\_Kurs\_Dane}
\begin{itemize}
  \item Kurs\_ID, Nazwa, ECTS, Semestr
  \item Data\_Wejscia, Data\_Wyjscia
\end{itemize}

\subsubsection*{SAT\_Oceny}
\begin{itemize}
  \item Student\_ID, Kurs\_ID, Data\_Egzaminu
  \item Ocena, Poprawka (boolean)
  \item Data\_Wejscia, Data\_Wyjscia
\end{itemize}

\newpage

\subsection{Schemat (opis tekstowy)}

\begin{verbatim}
HUB_Student           HUB_Prowadzacy
|                      |
SAT_Student_Dane       SAT_Prowadzacy_Dane
|                      |
|                 HUB_Kurs
|                      |
LINK_Student_Kurs    LINK_Kurs_Prowadzacy
|
SAT_Oceny

\end{verbatim}

\subsection{Potencjalne rozszerzenie modelu}

\subsubsection*{Nowy HUB: Grupa\_zajeciowa}
\begin{itemize}
  \item Grupa\_ID (np. ID grupy)
\end{itemize}

\subsubsection*{LINK\_Student\_Grupa}
\begin{itemize}
  \item Student\_ID, Grupa\_ID
\end{itemize}

\subsubsection*{LINK\_Prowadzacy\_Grupa}
\begin{itemize}
  \item Prowadzacy\_ID, Grupa\_ID
\end{itemize}

\subsubsection*{LINK\_Grupa\_Sala}
\begin{itemize}
  \item Grupa\_ID, Sala\_ID
\end{itemize}

\subsubsection*{SAT\_Grupa\_Dane}
\begin{itemize}
  \item Grupa\_ID, Typ\_zajec, Dzien\_tygodnia, Godzina
  \item Data\_Wejscia, Data\_Wyjscia
\end{itemize}

\subsection{Odpowiedzi na pytania dodatkowe}

\begin{itemize}
  \item \textbf{Gdzie są przechowywane zmiany (np. nazwiska)?} \\
        W satelitach \textit{SAT\_Student\_Dane}, \textit{SAT\_Prowadzacy\_Dane} - każda zmiana to nowy rekord z nową datą wejścia.

  \item \textbf{Gdzie są informacje o ocenie z egzaminu poprawkowego?} \\
        W \textit{SAT\_Oceny} - nowy rekord z tą samą kombinacją student-kurs, ale z nową datą i flagą \texttt{Poprawka = True}.

  \item \textbf{Jak oddzielić dane wrażliwe?} \\
        Umieścić je w osobnych satelitach (np. \textit{SAT\_Student\_Prywatne}) i zabezpieczyć dostęp na poziomie bazy.

  \item \textbf{Jak obsłużyć rozszerzenie modelu o grupy?} \\
        Dodać \textit{HUB\_Grupa\_zajeciowa} i powiązane linki - Data Vault umożliwia łatwe rozszerzanie bez naruszenia istniejącej struktury.
\end{itemize}

\section{Ocena modelu wygenerowanego przez sztuczną inteligencję}

ChatGPT nie poradził sobie najgorzej i nawet stworzył diagram przedstawiający relacje. Ale popełnił oczywiście różne błędy.

\end{document}
